\
\section{Application: epistemic anatomy of the IAEA Norcia case}
\label{sec:application}

The end-to-end workflow is demonstrated on the Norcia normal-faulting
case of the IAEA PFDHA exercise \citep[Section~7]{IAEA2025TECDOC2092}.
Section~\ref{sec:application:anchor}. The
substance of this section is what the framework adds on top of the
benchmark: the same two-fault source model is re-evaluated under the
\emph{complete} published model library applicable to \emph{tectonic}
normal faulting,
under the rupture-location weight implied by the mapping accuracy that
the IAEA exercise assigns to the fault traces
(Sections~\ref{sec:application:tunnel}--\ref{sec:application:anatomy}).

\subsection{Case set-up and verification anchor}
\label{sec:application:anchor}

The Norcia case combines the Mount Vettore fault (MVF; 34~km,
$47^\circ$~SW dip, characteristic Gaussian MFD centred on
$M_{\mathrm{w}} = 6.7$, mean rate $4.03 \times 10^{-4}$~a$^{-1}$) and
the Norcia fault (NF; 29~km, $50^\circ$~SW dip, truncated
Gutenberg--Richter with $b = 1$ up to $M_{\mathrm{w}} = 6.8$, mean rate
$5.41 \times 10^{-3}$~a$^{-1}$), both with normal kinematics
\citep[Table~14]{IAEA2025TECDOC2092}. Three sites from the IAEA site
set \citep[Table~13]{IAEA2025TECDOC2092} are evaluated, all three
sited on the Mount Vettore fault system: the principal site PF on the
southern MVF trace itself ($l/L = 0.05$); and two distributed sites on
antithetic structures that were reactivated in 2016 -- MS on the Mount
Serra fault, at the San Benedetto tunnel position ($r = 7.6$~km from
the MVF hanging wall, $r = 2.9$~km from the NF footwall), and SL on the
San Lorenzo fault at its closest approach to the MVF trace
($r = 2.4$~km, $l/L = 0.4$). With the two-chain logic tree
of the exercise (a Youngs et al. 2003 chain and a hybrid chain
replacing the distributed models with Visini et al. 2025), the
framework reproduces the IAEA reference curves at the three sites and
the regional map at the $10^{5}$~a return period; the corresponding
inputs and the run-by-run comparison are archived as a regression
benchmark in the repository. All IAEA sites carry the mapping-accuracy class ``Accurately located''
\citep[Table~13]{IAEA2025TECDOC2092}; following
\citet[Table~2]{Petersen2011}, this class corresponds to a two-sided
location uncertainty $\sigma = 26.89$~m, which is applied as the
rupture-location weight $W_p = \exp(-r^{2}/2\sigma^{2})$ on every
branch of every calculation below.

\subsection{What would a PFDHA have told the tunnel designer?}
\label{sec:application:tunnel}

Site MS is not an arbitrary grid point: it is the position where the
San Benedetto road tunnel was offset during the 2016
$M_{\mathrm{w}}$~6.6 Norcia mainshock, by 0.20~m vertical and 0.13~m
left-lateral slip (net ${\sim}0.24$~m) on the reactivated Mount Serra
antithetic fault \citep{Galli2019WTC, Galli2020}. This turns the site
into a concrete engineering question that a verification exercise
cannot ask: \emph{ex ante, what annual rate would a PFDHA have assigned
to the displacement that actually occurred?}

Answering it requires two pieces of bookkeeping first. The published
distributed-rupture library for normal faulting is not all in scope
here: \citet{FerrarioLivio2021} is calibrated on volcano-tectonic
normal faulting and is excluded, leaving the two models developed for
tectonic normal faulting, \citet{youngs_methodology_2003} and
\citet{Visini2025}. And distributed rupture probability is a
probability \emph{per cell}, with the cell belonging to the model
rather than to the site:
\citet{youngs_methodology_2003} is developed on a $500 \times 500$~m
cell and exposes no cell parameter -- the size is absorbed into the
fitted coefficients -- whereas \citet{Visini2025} publish coefficients
for cells of 10--500~m, selectable at run time. Visini is therefore
pinned to 500~m so that the two are comparable. The consequence
deserves stating plainly: the IAEA exercise specifies a
$100 \times 100$~m site \citep[Table~13]{IAEA2025TECDOC2092}, and no
published tectonic normal-faulting distributed model except
\citet{Visini2025} can answer at that scale.

Each model also carries its own published variants, which the tree
keeps as separate branches rather than averaging.
\citet{youngs_methodology_2003} provide two regressions of normalised
distributed displacement, fitted to the 85th and the 95th percentile of
the data scatter; \citet{Visini2025} depend on the
magnitude--area scaling relation used to derive the principal throw,
for which three published choices are carried. It is worth being
explicit about what each model delivers to the hazard integral, since
neither returns a point value. In \citet{youngs_methodology_2003} the
normalised displacement $D/MD$ follows a gamma distribution of shape
$a = 2.5$ whose scale is set so that its 85th (or 95th) percentile
equals the distance-dependent regression value, and this is then
integrated over a \citet{WellsCoppersmith1994} lognormal distribution
of the principal maximum displacement; the ``85th'' and ``95th'' labels
are thus alternative \emph{fits}, not alternative output percentiles.
In \citet{Visini2025} the regression predicts the \emph{median} throw,
with a lognormal residual ($\sigma_{\ln Y} = 1.03$) truncated at
$\pm 3\sigma$. Both therefore supply a full exceedance distribution,
but their central tendencies are anchored differently -- one at a high
percentile of normalised displacement, the other at the median of
absolute throw.

Figure~\ref{fig:ms_tunnel} shows the resulting curves at both
antithetic-fault sites. The dominant control turns out not to be the
choice of model at all, but a geological judgement about the site: is
there a pre-existing fault beneath it? In the \citet{Visini2025}
decision tree that question is the first branch, and answering it
changes the rate of exceeding the observed offset by a factor 137 at MS
and 49 at SL. By comparison, the choice between
\citeauthor{youngs_methodology_2003} and \citeauthor{Visini2025} spans
a factor 1.5, \citeauthor{youngs_methodology_2003}'s own
85th-versus-95th regression a factor 1.5, and \citeauthor{Visini2025}'s
three magnitude--area scaling relations a factor 1.0.

The 2016 observations discriminate between the two readings. Assuming
no pre-existing structure, the calculation returns return periods of
$2.8 \times 10^{6}$~a at MS and $7.0 \times 10^{6}$~a at SL for
displacements that were measured on the ground after the earthquake.
Admitting one, the same tree returns $2.0 \times 10^{4}$~a and
$1.3 \times 10^{5}$~a -- the right order for structures of this
activity. The evidence at these sites is unambiguous: MS and SL are
pre-existing antithetic faults of the Mount Vettore system and both
ruptured at the surface in 2016, displacing the San Benedetto tunnel at
MS and offsetting a road by of order 1~m at SL. What the comparison
demonstrates is that a site-characterisation decision, made before any
model is chosen, can outweigh the entire published model library by two
orders of magnitude.

This also sharpens a result that we had initially read more broadly.
Re-evaluating the identical \citet{Visini2025} chain at the 100~m IAEA
site dimension, changing nothing else, moves the rate of exceeding the
observed offset by a factor 8.2 when no pre-existing structure is
assumed, but by only 1.04 when one is. The cell convention is
consequential, but specifically for the ``simple'' distributed-rupture
pathway: only that pathway carries the along-strike Monte Carlo term,
whose value differs by a factor 6.6 between the two site dimensions.
Where reactivation of a mapped structure dominates, the along-strike
probability is unity by construction and the cell washes out. The
practical warning survives in narrowed form -- because
\citet{youngs_methodology_2003} exposes no cell parameter, a
cell-inconsistent chain is invisible in the input files -- and the
framework now records each model's calibration cell with its
implementation.

A single observation cannot validate any of these curves; what the
comparison shows is where two observations fall within the epistemic
range the library spans, once that range is measured on a consistent
basis and restricted to models in scope. One caveat remains: the
displacement definitions differ in component (vertical versus net), a
mismatch of order the 0.20/0.24 ratio here.

\subsection{The near-trace profile}
\label{sec:application:profile}

Because the mapping accuracy enters as a Gaussian weight truncated at
$2\sigma = 54$~m, the principal and distributed contributions hand over
to one another within a few tens of metres of the trace -- a transition
that no site-specific curve can show. We evaluate it directly on a
2-km profile perpendicular to the MVF trace through PF, sampled every
20~m, with the Youngs et al.\ (2003) chain fixed. The principal contribution exists only inside the $\pm 54$~m
band and peaks at 3.60~m at $10^{5}$~a; the corridor carrying
$D \geq 1$~m is ${\sim}110$~m wide, essentially the full
$\pm 2\sigma$ width (this profile's 20~m sampling only brackets it;
a 10~m sampling of the same profile would resolve it). Outside the band the total hazard
\emph{is} the distributed contribution, which falls off
asymmetrically: 0.49~m at $-100$~m on the hanging wall against 0.09~m
at $+100$~m on the footwall, and 0.41 against 0.06~m at $\pm 300$~m.
At the three IAEA sites this structure is invisible -- PF sits on the
trace, where $W_p = 1$ by construction, and MS and SL are kilometres
away -- which is the practical point: trace-location uncertainty is a
siting-scale parameter, first-order only within a few tens of metres of
the mapped trace, exactly the scale at which a fault-crossing lifeline
is designed.

\subsection{Epistemic anatomy}
\label{sec:application:anatomy}

Figure~\ref{fig:logic_tree} shows the tree in full: four
surface-rupture-probability models, three principal-displacement
branches, and the five in-scope distributed chains, giving 60 end
branches that are enumerated and aggregated in one run per job. Only
one level-to-level transition is a genuine pairing rather than a free
combination -- each distributed-rupture model brings its own published
displacement variants -- which is why the branch count is
$4 \times 3 \times (2+3) \times 1$ rather than a plain product over
five levels.

Per-branch curves at the three sites give the following. At PF the
$10^{5}$~a design displacement spans 2.04--3.58~m (weighted 5--95\%,
mean 2.53~m). At the two antithetic sites the weighted means
are 0.398~m at MS and 0.571~m at SL, with 95th percentiles of 0.66~m
and 0.92~m; the 5th percentile is zero at both, because a
majority-weight subset of branches places the $10^{5}$~a level below
the curve entirely.

The SL value is worth dwelling on. At 0.571~m it lies within a factor
1.8 of the ${\sim}1$~m offset observed there in 2016 -- a second check
on the site characterisation of
Section~\ref{sec:application:tunnel}, at a different site and a
different displacement level from the San Benedetto tunnel. Two
independent observations, at two sites, are consistent with the same
classification.

These distributed values are also a reminder that model applicability
screening is a first-order decision at such sites rather than
bookkeeping. An earlier version of this tree included
\citet{FerrarioLivio2021} and returned 1.8~cm at MS, supplied almost
entirely by that single volcano-tectonic branch. That figure was wrong
twice over: the model does not apply to these purely tectonic sources,
and the value was an order of magnitude below what the in-scope
library gives once the site is correctly characterised. A framework
that makes the model set an explicit, versioned input is what allows
either error to be seen and audited.

Two framework-level observations complete the picture. Kuehn et
al.\ (2024) could not legally enter this tree: it predicts
\emph{aggregate} displacement, and the validator rejects any chain that
combines an aggregate primary model with distributed models, because no
conversion between displacement definitions exists (Sarmiento et
al.\ 2025). Likewise, holding the Youngs (principal) and Lavrentiadis
\& Abrahamson (sum-of-principal) definitions in one tree required
splitting the primary level into definition-consistent branch sets --
visible in the input file itself. The epistemic menu is thus
constrained by displacement-definition bookkeeping, and the framework
machine-checks constraints that in current practice are enforced, if at
all, by footnote.

\subsection{Regional view}
\label{sec:application:map}

Figure~\ref{fig:anatomy_maps} maps the same tree at the $10^{5}$~a
return period: panel (a) the mean hazard over all 60 branches, panel
(b) the 84th epistemic fractile. A resolution caveat governs how such a
map may be read: the grid spacing, $0.01^{\circ}$ or 1112~m, is twenty
times wider
than the ${\sim}110$~m principal corridor, so whether a pixel falls
inside the $\pm 2\sigma$ band is an accident of the grid. Refining the
regional grid is not a practical remedy -- cost is linear in the number
of sites, so $0.001^{\circ}$ would take some seven hours and still
place barely one cell across the corridor.

We therefore restrict the map to the principal component, evaluated on
the exact on-trace sites that the engine inserts alongside the filtered
grid. This is not a presentational convenience but a consequence of
Section~\ref{sec:application:tunnel}: the distributed term depends on
whether a pre-existing structure lies beneath the site, which is a
per-site geological judgement -- case~1 at MS and SL, which sit on such
structures -- and applying it uniformly across a grid would
over-predict everywhere away from the antithetic faults. The principal
term carries no such dependence and is well posed at every point. It is
also intrinsically confined: of 992 grid sites only 8 fall within the
$\pm 2\sigma$ band, against all 80 on-trace sites. Near-trace structure
is characterised instead by the cross-trace profile of
Section~\ref{sec:application:profile}.

Panel~(c) of Figure~\ref{fig:anatomy_maps} shows where the mapped
values come from. At PF the mean hazard gives 2.53~m at $10^{5}$~a
against a 16th--84th epistemic band of 2.19--3.30~m, and the branches
separate mainly above 0.1~m: the low-displacement plateau is set by the
probability of surface rupture, on which the models largely agree at
$M_w$ 6.4--7.0, while the spread at design level is carried by the
principal displacement models. Reducing epistemic uncertainty at this
site would therefore mean constraining the displacement model, not the
rupture-probability model. Along strike the mean is 2.91~m and the 84th
fractile 3.76~m (medians over the on-trace sites), a 29\% difference
that is visible directly as the colour offset between panels (a) and
(b).

% ---------------------------------------------------------------------
% Figure environments (paths relative to the paper's Figures/ directory;
% copy from openquake/fdha/test/usercase/norcia_anatomy/figures/)
% ---------------------------------------------------------------------
\subsection{From hazard curve to design displacement}
\label{sec:application:design}

The maps report the displacement at a fixed return period read from the
mean hazard curve -- the equal-hazard convention. \citet{AbrahamsonLiou2025}
show that this convention can misrepresent the epistemic centre when
the displacement hazard curve is flat at the design return period, as
it is for moderate-activity faults, and propose alternatives: an upper
fractile of the hazard, a risk-targeted displacement, or the median
displacement given surface rupture at the site. Whether those
alternatives are needed is decided by the slope of the hazard curve,
which we therefore compute rather than assume.

For the on-fault principal hazard here the log--log slope at
$10^{5}$~a is $-2.40$, steeper than the
high-activity case of that paper ($-2.07$) and far from its
moderate-activity case ($-0.10$). The equal-hazard value therefore
stands: 2.91~m as the median over the on-trace sites. For completeness
we also evaluate their equation~(9), the distribution of displacement
conditional on rupture occurring, which returns a median of 0.20~m and
an 85th percentile of 0.91~m at these sites; the surface-rupture
recurrence is $3.4 \times 10^{3}$~a, so the $10^{5}$~a level sits far
out on that conditional tail. These conditional values are reported as
diagnostics only. \citeauthor{AbrahamsonLiou2025} are explicit that the
approach is intended for flat hazard curves and is not to be used to
reduce the displacement where the curve is steep, and substituting
0.91~m for 2.91~m here would be exactly that misuse. Their closing
recommendation -- that PFDHA reports should present the epistemic
fractiles of the displacement hazard, as modern PSHA reports do -- is
met by Figure~\ref{fig:anatomy_maps}.

\begin{figure}[!htbp]
  \centering
  \includegraphics[width=0.98\linewidth]{fig_logic_tree.pdf}
  \caption{The FDHA logic tree: five levels, $4 \times 3 \times (2+3)
  \times 1 = 60$ end branches. Weights are normalised, and the product
  along any path gives that branch's weight. Links are drawn only where
  the levels do not combine freely -- each distributed-rupture model
  carries its own published displacement variants, so the
  distributed-rupture level is subsumed in the pairing term of that
  count. The $0.25/0.25/0.5$ weighting of the principal-displacement
  level is what expresses the two displacement definitions (Youngs et
  al.\ principal, Lavrentiadis \& Abrahamson sum-of-principal) being
  held in separate definition-consistent branch sets, a constraint the
  framework enforces rather than leaves to a footnote.}
  \label{fig:logic_tree}
\end{figure}

\begin{figure}[!htbp]
  \centering
  \includegraphics[width=0.95\linewidth]{fig_ms_tunnel.pdf}
  \caption{Mean distributed-displacement hazard at the two
  antithetic-fault sites, (a)~MS and (b)~SL, against what was measured
  at each in 2016 (vertical dashed lines: the 0.20~m vertical offset of
  the San Benedetto tunnel, and a road offset of order 1~m at San
  Lorenzo). Each curve is the weighted mean over the primary nodes that
  gate it. Blue: the two \citet{youngs_methodology_2003} regressions,
  fitted to the 85th and 95th percentile of the data scatter. Red and
  orange: \citet{Visini2025} evaluated with and without a pre-existing
  fault assumed beneath the site, the first branch of that paper's site
  decision tree. Solid curves use the $500 \times 500$~m cell, the
  fixed calibration cell of \citet{youngs_methodology_2003}, matched by
  pinning the \citet{Visini2025} \texttt{pixel\_size}; dashed curves are
  the identical chain at the 100~m IAEA site dimension and differ
  \emph{only} in that convention. The cell separates the curves when no
  pre-existing fault is assumed and not otherwise, because only the
  simple-rupture pathway carries the cell-dependent along-strike term.}
  \label{fig:ms_tunnel}
\end{figure}

\begin{figure}[!htbp]
  \centering
  \includegraphics[width=0.95\linewidth]{fig_principal_fractiles.pdf}
  \caption{Principal displacement at the $10^{5}$~a return period over
  the 60-branch tree: (a)~mean hazard, (b)~84th epistemic fractile,
  both on the on-fault sites only; (c)~every branch at PF (grey) with
  the mean hazard (black) and the 16th--84th epistemic band (shaded).
  Triangles: the three IAEA sites. Only the principal component is
  mapped: the distributed term depends on the per-site
  characterisation of Section~\ref{sec:application:tunnel} and cannot
  be applied uniformly over a grid, whereas the principal term carries
  no such dependence. It is non-zero only within the $\pm 2\sigma$
  rupture-location band, that is, on the fault. The fractiles are taken
  across logic-tree branches and are therefore epistemic; aleatory
  variability is already integrated within each individual branch
  curve.}
  \label{fig:anatomy_maps}
\end{figure}

